\documentclass[11pt]{article}
\usepackage{amsmath,amssymb,amsthm,mathtools}
\usepackage{enumitem}
\usepackage{hyperref}
\title{Finite-Window Kummer Recognition of the Canonical Orientation\\of a Rank-Four Demu\v{s}kin Pro-$3$ Group}
\author{}
\date{September 2026}
\newtheorem{theorem}{Theorem}[section]
\newtheorem{proposition}[theorem]{Proposition}
\newtheorem{lemma}[theorem]{Lemma}
\begin{document}
\maketitle

\begin{abstract}
We study finite-window recognition of the canonical orientation of the fixed rank-four q=3 Demuskin pro-3 group. Let $P_i$ be the Zassenhaus filtration and put $Q_k=G/P_{k+1}$. For a principal-unit character $\rho:Q_k\to U_{1,k}:=1+3(\mathbb Z/3^k\mathbb Z)$, let $\mathsf K_k(Q_k,\rho)$ mean that $H^1(Q_k,\mathbb Z/3^k(\rho))\to H^1(Q_k,\mathbb F_3)$ is surjective. We prove that $\mathsf K_k(Q_k,\rho)$ holds exactly for $\rho=\chi_G\bmod3^k$, for $k\ge2$. The canonical orientation itself is classical; the result is a finite-window recognition statement.
\end{abstract}

\section{Introduction}
Let
\[
G=\langle x_1,x_2,x_3,x_4 \mid x_1^3[x_1,x_2][x_3,x_4]=1\rangle.
\]
An orientation of a pro-$3$ group $H$ is a continuous homomorphism $\theta:H\to1+3\mathbb Z_3$, giving the twisted module $\mathbb Z_3(\theta)$. The oriented group is Kummerian when, for every n>=1, the natural map H^1(H,Z_3(theta)/3^n) to H^1(H,F_3) is surjective; this is the standard Kummerian formulation used by Efrat--Quadrelli and Quadrelli--Weigel \cite{EfratQuadrelli2019,QuadrelliWeigel2020,QuadrelliWeigel2022}. Labute proved that every Demushkin group has a unique orientation making it Kummerian, namely the canonical orientation induced by the dualizing module \cite{Labute1967}. Our question is whether its reduction modulo 3^k can be recognized from Q_k=G/P_{k+1} without supplying the orientation as input.

\begin{theorem}[Finite-window recognition]
For every k>=2 and every principal-unit rho:Q_k to U_{1,k}:=1+3A_k,
\[
\mathsf K_k(Q_k,rho) \Longleftrightarrow rho=chi_G\pmod {3^k}.
\]
\end{theorem}

The proof separates five logically distinct steps: the finite semidirect-product filtration, arbitrary-candidate factorization, the one-relator lifting obstruction, the base-level identification, and intrinsic PD$^2$ uniqueness.

\section{Finite window and arbitrary-candidate factorization}
Put $A_k=\mathbf Z/3^k\mathbf Z$, $U_j=1+3^jA_k$ for $j\ge1$, with $U_j=1$ for $j\ge k$, and
\[
S_k=A_k\rtimes U_1,\qquad (a,u)(b,v)=(a+ub,uv).
\]
Let $P_1(S_k)=S_k$ and $P_{j+1}(S_k)=P_j(S_k)^3[P_j(S_k),S_k]$.

\begin{lemma}[U1: finite-depth semidirect filtration]
For $1\le j\le k+1$, put $T_j=3^{j-1}A_k\rtimes U_j$. Then
\[
P_j(S_k)=T_j,\qquad P_{k+1}(S_k)=1.
\]
\end{lemma}

\begin{proof}
For $(a,u)\in T_j$, write $u=1+3^jt$. Then
\[
(a,u)^3=((1+u+u^2)a,u^3),
\]
and $1+u+u^2\in3A_k$. Moreover $U_j^3=U_{j+1}$ for $j<k$ by successive $3$-adic digit lifting. Hence $T_j^3\subseteq3^jA_k\rtimes U_{j+1}$.

A direct commutator calculation gives $[T_j,S_k]\subseteq3^jA_k$. Conversely, with $[x,y]=x^{-1}y^{-1}xy$ and $4\in U_1$,
\[
[(a,1),(0,4)]=((4^{-1}-1)a,1).
\]
Since $4^{-1}-1=-3/4$ is $3$ times a unit, these commutators generate all of $3^jA_k$. Thus $[T_j,S_k]=3^jA_k$. Induction from $P_1=T_1$ now gives
\[
P_{j+1}(S_k)=T_j^3[T_j,S_k]=3^jA_k\rtimes U_{j+1}=T_{j+1}.
\]
At $j=k$, both factors are trivial.
\end{proof}

\begin{lemma}[U2: arbitrary-candidate factorization]
Let $G$ be any pro-$3$ group, $\rho:G\to U_{1,k}$ a character, and
$z\in Z^1(G,A_k(\rho))$. Then $z$ and $\rho$ factor through
$Q_k=G/P_{k+1}(G)$, and inflation gives a canonical isomorphism
\[
H^1(Q_k,A_k(\bar\rho))\xrightarrow{\sim}H^1(G,A_k(\rho)).
\]
\end{lemma}

\begin{proof}
The crossed-cocycle identity is equivalent to $\psi(g)=(z(g),\rho(g))$ being a homomorphism $G\to S_k$. Functoriality of the lower $3$-central series and U1 give $\psi(P_{k+1}(G))=1$. Thus both components factor through $Q_k$. The same argument for coboundaries identifies $B^1$ on both sides, proving the cohomology isomorphism.
\end{proof}

\section{Finite Kummer criterion for a minimal one-relator presentation}
Let $F$ be free pro-$3$ on $x_1,\dots,x_d$, let $G=F/\overline{\langle\!\langle r\rangle\!\rangle}$, and assume $r\in\Phi(F)$. For a principal-unit character $\rho:G\to U_{1,k}$, write $F_i(\rho)\in A_k$ for the twisted Fox evaluation, with the convention
\[
z(r)=\sum_iF_i(\rho)z(x_i).
\]

\begin{proposition}[U3: finite Kummer/Fox criterion]
Under this minimality hypothesis,
\[
\mathsf K_k(G,\rho)\Longleftrightarrow F_i(\rho)=0
\quad(1\le i\le d).
\end{proposition}

\begin{proof}
Minimality gives $H^1(G,\mathbf F_3)\cong\mathbf F_3^d$ by generator values. A vector $a=(a_i)\in A_k^d$ is the value vector of a crossed cocycle exactly when $\sum_iF_i(\rho)a_i=0$.

Assume every mod-$3$ vector lifts. For the lift of the $i$-th basis vector, write $a_i=1+3b_i$ and $a_j=3b_j$ for $j\ne i$. The relation equation first gives $F_i\in3A_k$. If all $F_j\in3^mA_k$, the same equation gives
\[
0=F_i(1+3b_i)+\sum_{j\ne i}F_j(3b_j)
=F_i+3\sum_jF_jb_j,
\]
so $F_i\in3^{m+1}A_k$. Induction from $m=0$ to $m=k$ gives $F_i=0$. Conversely, if all $F_i=0$, every generator-value vector in $A_k^d$ satisfies the relation, so every mod-$3$ class lifts.
\end{proof}

The principal-unit restriction makes the reduction of the coefficient action trivial modulo $3$. The argument is valuation induction, not the earlier informal Nakayama step.

\section{Coordinate identification at the base level}
For
\[
r=x_1^3[x_1,x_2][x_3,x_4],
\]
the standard odd-$p$ Demu\v{s}kin calculation gives
\[
rho(x_1)=rho(x_3)=rho(x_4)=1, \qquad
rho(x_2)=(1-3)^{-1}\pmod {3^k}.
\]
Thus the values at $k=2,3,4$ are $4\bmod9$, $13\bmod27$, and $40\bmod81$.

\section{Intrinsic uniqueness}
The next three results are the load-bearing U5 argument.

\begin{lemma}[U5a: reduction]
If $\mathsf K_k(G,\rho_k)$ holds, then $\mathsf K_{k-1}(G,\rho_{k-1})$ holds after reduction modulo $3^{k-1}$.
\end{lemma}

\begin{proof}
Reduce a crossed cocycle modulo $3^{k-1}$. A level-$k$ lift of any class in $H^1(G,\mathbf F_3)$ therefore gives a level-$(k-1)$ lift of the same class.
\end{proof}

\begin{lemma}[U5b: coefficient-extension variation]
Suppose $\rho_k$ and $\rho_k'$ lift the same $\rho_{k-1}$ and
\[
\rho_k'=\rho_k(1+3^{k-1}\nu),\qquad \nu\in Z^1(G,\mathbf F_3).
\]
Let $\delta_{\rho_k}$ and $\delta_{\rho_k'}$ be the connecting maps of the corresponding coefficient extensions. Then
\[
\boxed{\delta_{\rho_k'}-\delta_{\rho_k}=\iota_*\circ(\nu\smile -),}
\]
where $\iota_*:H^2(G,\mathbf F_3)\to H^2(G,A_{k-1}(\rho_{k-1}))$ is induced by $\iota(1)=3^{k-2}$.
\end{lemma}

\begin{proof}
Choose the common set-theoretic section $s:\mathbf F_3\to A_k$ with $s(1)=1$. For $c\in Z^1(G,\mathbf F_3)$, the connecting class is represented by the defect
\[
g s(c(h))+s(c(g))-s(c(gh)).
\]
The two coefficient actions agree modulo $3^{k-1}$. Their difference on the lifted value $s(c(h))$ is the socle term $3^{k-1}\nu(g)c(h)$. Under the identification of the socle of the kernel with $\mathbf F_3$ via $1\mapsto3^{k-2}$, this additional defect is precisely the cochain representing $\iota(\nu\smile c)$. The cocycle identities show that the difference is a $2$-cocycle, and changing the section changes it by a coboundary. Therefore the displayed identity holds in cohomology.
\end{proof}

\begin{lemma}[U5c: PD$^2$ socle injectivity]
For the canonical orientation $\chi$, the inclusion
\[
\iota:\mathbf F_3\hookrightarrow A_{k-1}(\chi_{k-1}),\qquad
1\mapsto3^{k-2},
\]
induces an injection
\[
\iota_*:H^2(G,\mathbf F_3)\hookrightarrow
H^2(G,A_{k-1}(\chi_{k-1})).
\]
\end{lemma}

\begin{proof}
Let $D=\mathbf Q_3/\mathbf Z_3(\chi)$ be the discrete dualizing module. Finite-module PD$^2$ duality gives
\[
H^2(G,M)^\vee\cong\operatorname{Hom}_G(M,D).
\]
For $M=A_{k-1}(\chi_{k-1})$ this Hom group is cyclic of order $3^{k-1}$, while for $M=\mathbf F_3$ it is $\mathbf F_3$, since $D[3]\cong\mathbf F_3$ with trivial residual action. Under these identifications, precomposition with $\iota$ is the reduction map
$\mathbf Z/3^{k-1}\twoheadrightarrow\mathbf F_3$. Thus the dual of $\iota_*$ is surjective and $\iota_*$ is injective.
\end{proof}

\begin{proposition}[U5: intrinsic uniqueness]
If $\rho_k$ and $\rho_k'$ both satisfy $\mathsf K_k$ and both reduce to $\chi_{k-1}$, then $\rho_k=\rho_k'$.
\end{proposition}

\begin{proof}
The level-$k$ predicate is equivalent to vanishing of the connecting map. Hence U5b gives
\[
0=\delta_{\rho_k'}-\delta_{\rho_k}=\iota_*(\nu\smile -).
\]
By U5c, $\nu\smile v=0$ for every $v\in H^1(G,\mathbf F_3)$. The Demu\v{s}kin cup pairing is nondegenerate, hence $\nu=0$.
\end{proof}

\section{Proof of the theorem}
Existence of chi_G with the Kummerian lifting property is classical. Factorization through Q_k was proved above. For uniqueness, reduce a candidate successively to level 2. The standard calculation gives rho_2=chi_G mod 9. Inductively, the candidate and chi_G mod 3^k have the same reduction and both satisfy the predicate; intrinsic uniqueness gives equality.

\section{Literature boundary and contribution}
The existence and uniqueness of the canonical Demu\v{s}kin orientation, and the corresponding all-level Kummerian criterion, are classical. Labute gives the canonical orientation in the classification of Demu\v{s}kin groups \cite{Labute1967}; the cohomological and cocycle formulations of Kummerianity are developed in Efrat--Quadrelli and subsequent work \cite{EfratQuadrelli2019,QuadrelliWeigel2020,QuadrelliWeigel2022}. We therefore do not claim the canonical orientation, or the global Kummerian criterion, as new.

The distinction relevant here is the direction of recognition. The classical theory starts with an oriented pair $(G,\theta)$ and asks whether the prescribed orientation is Kummerian. In contrast, our finite-window statement starts from the bare finite quotient
\[
Q_k=G/P_{k+1}
\]
and an arbitrary principal-unit candidate $\rho:Q_k\to U_{1,k}$, and uses the intrinsic predicate $\mathsf K_k(Q_k,\rho)$ to select the reduction of the canonical orientation. The finite-depth factorization in Sections 2--3 is part of this direction: the candidate is not supplied by the prior theory, and the relevant twisted cocycles are shown to factor through $Q_k$ for an arbitrary candidate.

A potentially relevant quotient-inheritance result is Proposition 2.10 of Quadrelli \cite{Quadrelli2024}. It assumes an already Kummerian oriented pair $(G,\theta)$, a normal subgroup $N\subseteq\ker(\theta)$, and surjectivity of
\[
H^1(G,\mathbb F_p)\longrightarrow H^1(N,\mathbb F_p)^G.
\]
Those hypotheses are materially different from the present recognition problem. In the present rank-four $p=3$ case, the natural choice $N=P_3$ does not satisfy the required restriction-surjectivity condition: $P_3\subseteq\Phi(G)$, whereas $P_3/P_4\ne0$ (indeed its dimension is $20$). Thus Proposition 2.10 cannot be invoked as an automatic quotient-inheritance mechanism for our finite-window theorem. This is a separation observation, not a claim that Proposition 2.10 itself is new or that no other special quotient could satisfy its hypotheses.

Within the literature audited for this manuscript, we did not identify an exact theorem whose input is the bare $Q_k$, whose candidate $\rho$ is arbitrary, whose selector is precisely the finite Kummer predicate $\mathsf K_k(Q_k,\rho)$, and whose arbitrary-candidate factorization through $Q_k$ is established. Accordingly, the publication-level novelty claim is intentionally conditional and narrow; it is not a claim of absolute priority.

\section{Limitations and logical boundaries}
The theorem is for the fixed rank-four $q=3$ group. ``$q$-blind'' means that $q$ is absent from the selector input; it does not mean uniformity over the Demu\v{s}kin family. $P_{k+1}$ is sufficient; minimality of the window is open. The final predicate is intrinsic although the standard presentation is used as a local calculation.

\section{Conclusion}
The result separates a classical infinite-group existence theorem from a finite recognition principle. It should therefore be read as a theorem about the information required by a selector, not as a new construction of the canonical orientation:
\[
G/P_{k+1} \quad+\quad \text{intrinsic Kummer lifting}
\quad\Longrightarrow\quad \chi_G\pmod {3^k}.
\]

\bibliographystyle{amsplain}
\bibliography{references}
\end{document}